\section{Discussion and Limitations}
\label{sec:discussion}

\subsection{Stable storage is the reusable interface}

Dynamic identity and graph replay conflict only when identity changes alter
graph-visible allocation or execution structure. A fixed storage interface
moves dynamism into contents and metadata updated at a controlled boundary.
This principle may apply to other bounded working sets with a legal staging
boundary and enforceable publication/reuse order; our evidence covers only the
implemented MoE path.

Address stability does not prevent stale publication or eviction; owner/version
checks and stream dependencies enforce temporal correctness. Static residency
sacrifices bounded offload, pointer rebinding requires graph recapture, and
eager fallback pays dispatch overhead. LatchMoE fixes slots and stages
replacement; its seam and fused operator are hardware-specific.

\subsection{Control policy is outside the data plane}

LatchMoE accepts an ordered placement plan but constrains it to the actual
routed expert set. This boundary lets future caching or prediction policies
reuse the fixed slots, transfer engine, wave executor, and graph checks. The
current evaluation does not compare placement policies and should not be read
as evidence that the default ordering is optimal. Such a study requires a
shared trace, identical storage capacity, and policy-specific hit, transfer,
and latency measurements on the same data plane.

\subsection{Scope of the current evidence}

The semantic and graph qualification covers three unquantized BF16 models on
one Ascend 910B2 with tensor parallelism one. It validates native router and
layer-boundary parity, exact end-to-end tokens, replay-visible addresses,
overflow composition, lease/release records, and zero eager fallback for the
three reported capability tuples. Performance experiments are narrower:
Qwen3-30B-A3B, a 14-GiB host-offload target, serialized requests,
\texttt{max\_num\_seqs=1}, and disabled prefix caching. The matched result
shows lower TTFT for multi-wave than full-layer staging under that contract;
the one-factor campaign isolates an overlap contribution. The full-resident
anchor also exposes the absolute cost of host-backed replacement. Native
prefetch and legacy layered offload complete graph-mode requests but fail the
exact-token comparability gate, so we retain them without using their latency.
Consequently, the reported performance evidence is an internal mechanism
comparison against graph-compatible full-layer staging, not a superiority claim
over published offloading systems.
The slot sweep has one start per point and is descriptive. None of these
results establishes a general decode-speedup or multi-request serving claim.

Higher concurrency permits overlapping request leases, prefix-cache reuse adds
hits across the staging seam, multi-device execution adds distributed
ownership, and quantization changes layouts and operator support. None is
covered by the qualified single-request host-to-device protocol.
